% !PRINTED 218
\weberpage{218}

Es folgt nämlich zunächst aus der zweiten der Relationen (11):
\[
\Theta^2\omega = \Theta\omega\Theta^3 = \omega\Theta^2, \quad \Theta^3\omega = \Theta\omega\Theta^2 = \omega\Theta,
\]
ferner aus der letzten (11):
\[
\Theta^3\chi = \Theta\chi\omega\Theta^3 = \chi^2\omega\Theta^2\omega\Theta^3 = \chi^2\Theta,
\]
und weiter:
\[
\omega\chi^2 = \chi^2\omega\chi = \chi\omega,
\]
\[
\Theta^2\chi^2 = \chi\omega\Theta^3\chi = \chi\omega\chi^2\Theta = \chi^2\omega\Theta.
\]
Wenn wir nun irgend drei Elemente $\chi$, $\omega$, $\Theta$ haben, die
sich nach irgend einer Regel componiren lassen, wenn dabei $\chi$ vom
dritten, $\omega$ vom zweiten, $\Theta$ vom vierten Grade ist, so folgt aus
dem Bestehen der Relationen (11), dass die Elemente
$\chi^{\lambda}\omega^{\mu}\Theta^{\nu}$ eine Gruppe 24\textsuperscript{sten} Grades bilden, die mit
der Octaëdergruppe isomorph ist. Dazu ist nur noch nachzuweisen,
dass aus diesen Voraussetzungen folgt, dass die 24 Elemente
$\chi^{\lambda}\omega^{\mu}\Theta^{\nu}$ alle von einander verschieden sind. Nehmen wir
an, es sei
\[
\chi^{\lambda}\omega^{\mu}\Theta^{\nu} = \chi^{\lambda'}\omega^{\mu'}\Theta^{\nu'},
\]
so würde folgen:
\[
\chi^{\lambda - \lambda'} = \omega^{\mu'}\Theta^{\nu' - \nu}\omega^{-\mu},
\]
und nach (11):
\[
\chi^{\lambda - \lambda'} = \omega^{\mu' - \mu}\Theta^{\pm(\nu - \nu')}.
\]
Nun folgt aber aus (11), dass $\omega\Theta$, $\omega\Theta^2$, $\omega\Theta^3$ vom
zweiten Grade sind, und da $\chi$ vom dritten Grade ist, so kann diese
Beziehung nur stattfinden, wenn $\lambda - \lambda'$ durch 3, $\mu - \mu'$ durch 2 und
$\nu - \nu'$ durch 4 theilbar, also $\chi^{\lambda} = \chi^{\lambda'}$, $\omega^{\mu} = \omega^{\mu'}$, $\Theta^{\nu} = \Theta^{\nu'}$ ist.

Nach der aus (10) folgenden Relation
\[
\omega = \chi\Theta\chi\Theta^2
\]
können wir $\omega$ aus $\chi$ und $\Theta$ zusammensetzen und daher $\chi$ und
$\Theta$ als \gesperrt{erzeugende Substitutionen} der Gruppe ansehen.

Die noch fehlenden beiden Grundformen achten und zwölften
Grades findet man sehr leicht, wenn man zunächst die Hesse'sche
Covariante von $f$ bildet.

Man erhält so die Grundform achten Grades:
\weq{13}{W = x_1^8 + 14\,x_1^4x_2^4 + x_2^8,}
und wenn man aus $f$ und $W$ die Functionaldeterminante bildet,
die ja wieder eine Covariante von $f$ ist (Bd. I, \S.\,59, 60), so
ergiebt sich die Grundform zwölften Grades:
\weq{14}{K = x_1^{12} - 33\,x_1^8x_2^4 - 33\,x_1^4x_2^8 + x_2^{12}.}
